%! Trigonometry and Polar Coordinates — Unit 7, Lesson 7.6 — Group Activity (Answer Key)
\documentclass[10pt]{article}
\usepackage{precalculus-article}
\usepackage{precalculus-key}   % pulls in -boxes; do NOT also load -boxes
\newcommand{\TallMath}[1]{$\displaystyle #1\rule[-1.4em]{0pt}{3.2em}$}

\begin{document}

\pageheader{Unit 7, Lesson 7.6}{Group Activity --- Answer Key}
\namepartnerperiod

\vspace{0.10in}

%----------------------------------------------------------------------
% Context
%----------------------------------------------------------------------
\begin{scenariobox}[Context: Polar Coordinates in Navigation and Radar]{plumlight}
Air traffic control radar systems report the location of aircraft using a
\emph{distance} (range) and a \emph{bearing} (angle from north) --- not $x$
and $y$ coordinates. This is the polar coordinate system in action. In this
activity your group will practice converting between polar and rectangular
representations and explore the fact that a single location has many valid
polar descriptions.
\end{scenariobox}

\vspace{0.08in}

%----------------------------------------------------------------------
% Tier R
%----------------------------------------------------------------------
\begin{tcolorbox}[colback=white, colframe=black!40,
    title=\textbf{Tier R --- Reinforce},
    fonttitle=\bfseries, arc=2mm, left=3mm, right=3mm, top=2mm, bottom=2mm]

\textbf{Directions:} For each polar coordinate pair, find the rectangular
coordinates using $x = r\cos\theta$ and $y = r\sin\theta$.

\vspace{6pt}
\renewcommand{\arraystretch}{2.2}
\begin{tabular}{|c|c|c|c|c|}
\hline
\textbf{Polar point} & \textbf{$r$} & \textbf{$\theta$} & \textbf{$x = r\cos\theta$} & \textbf{$y = r\sin\theta$} \\
\hline
$(4,\, \pi/6)$ & 4 & $\pi/6$ & \ans{$4 \cdot \frac{\sqrt{3}}{2} = 2\sqrt{3}$} & \ans{$4 \cdot \frac{1}{2} = 2$} \\
\hline
$(2,\, \pi/2)$ & 2 & $\pi/2$ & \ans{$2 \cdot 0 = 0$} & \ans{$2 \cdot 1 = 2$} \\
\hline
$(6,\, \pi)$ & 6 & $\pi$ & \ans{$6 \cdot (-1) = -6$} & \ans{$6 \cdot 0 = 0$} \\
\hline
$(3,\, 3\pi/2)$ & 3 & $3\pi/2$ & \ans{$3 \cdot 0 = 0$} & \ans{$3 \cdot (-1) = -3$} \\
\hline
\end{tabular}

\vspace{8pt}
\textbf{1.}\quad
$\sqrt{(2\sqrt{3})^2 + 2^2} = \sqrt{\ans{12} + \ans{4}} = \sqrt{\ans{16}} = $ \ans{$4$} \quad
Does this equal the original $r$? \ans{Yes}

\vspace{8pt}
\textbf{2.}\quad \ans{The angle $\pi/2$ points along the positive $y$-axis. The horizontal
component of any point on the $y$-axis is zero, so $x = 0$.}

\vspace{6pt}
\textbf{3.}\quad Rectangular: \ans{$(5,\, 0)$} \quad Reason: \ans{$\theta = 0$ means the point is on the positive $x$-axis, so $y = 0$ and $x = r = 5$.}

\end{tcolorbox}

\vspace{0.08in}

%----------------------------------------------------------------------
% Tier A
%----------------------------------------------------------------------
\begin{tcolorbox}[colback=white, colframe=black!40,
    title=\textbf{Tier A --- Apply},
    fonttitle=\bfseries, arc=2mm, left=3mm, right=3mm, top=2mm, bottom=2mm]

\textbf{Part 1 --- Rectangular $\to$ Polar}

\vspace{4pt}
(a)\quad $(1,\, \sqrt{3})$: \quad $r = $ \ans{$2$} \quad
$\theta = $ \ans{$\pi/3$} \quad Polar: \ans{$(2,\; \pi/3)$}

\vspace{6pt}
(b)\quad $(-4,\, 0)$: \quad $r = $ \ans{$4$} \quad
$\theta = $ \ans{$\pi$} \quad Polar: \ans{$(4,\; \pi)$}

\vspace{6pt}
(c)\quad $(3,\, -3)$: \quad $r = $ \ans{$3\sqrt{2}$} \quad
$\theta = $ \ans{$7\pi/4$} \quad Polar: \ans{$(3\sqrt{2},\; 7\pi/4)$}

\vspace{10pt}
\textbf{Part 2 --- Multiple Representations}

\vspace{4pt}
Representation 1: \ans{$(2,\; \pi/3 + 2\pi) = (2,\; 7\pi/3)$}

\vspace{4pt}
Representation 2: \ans{$(-2,\; \pi/3 + \pi) = (-2,\; 4\pi/3)$}

\vspace{4pt}
$(2,\pi/3)$ in rectangular: \ans{$(1,\; \sqrt{3})$}
\quad Rep.\ 2 in rectangular: $(-2)\cos(4\pi/3) = (-2)(-1/2) = 1$; $(-2)\sin(4\pi/3) = (-2)(-\sqrt{3}/2) = \sqrt{3}$; \ans{$(1,\; \sqrt{3})$}
\quad Match? \ans{Yes}

\vspace{10pt}
\textbf{Part 3 --- Complex Number in Polar Form}

\vspace{4pt}
$r = \sqrt{(\ans{2})^2 + (\ans{-2})^2} = $ \ans{$2\sqrt{2}$}

\vspace{4pt}
Quadrant of $2-2i$: \ans{Q4} \qquad
$\theta = \arctan\!\left(\dfrac{-2}{2}\right) = $ \ans{$-\pi/4$}
(Q4, $x > 0$, no adjustment needed; use $7\pi/4$): $\theta = $ \ans{$7\pi/4$}

\vspace{4pt}
Polar form: \ans{$2\sqrt{2}\cos(7\pi/4) + 2\sqrt{2}\,i\sin(7\pi/4)$}

\end{tcolorbox}

\vspace{0.08in}

%----------------------------------------------------------------------
% Tier E
%----------------------------------------------------------------------
\begin{tcolorbox}[colback=white, colframe=black!40,
    title=\textbf{Tier E --- Extend},
    fonttitle=\bfseries, arc=2mm, left=3mm, right=3mm, top=2mm, bottom=2mm]

\textbf{Problem 1}

\vspace{4pt}
(a)\quad $(r,\theta) = $ \ans{$(r,\; \pi - \alpha)$}

\vspace{6pt}
(b)\quad All representations: \ans{$(r,\; \pi - \alpha + 2\pi k)$} for any integer $k$,
or \ans{$(-r,\; -\alpha + 2\pi k)$} for any integer $k$.

\vspace{8pt}
\textbf{Problem 2}

\vspace{4pt}
Rectangular coordinates of $(-r,\theta)$:
\begin{align*}
x_1 &= (-r)\cos\theta = \ans{$-r\cos\theta$} \\
y_1 &= (-r)\sin\theta = \ans{$-r\sin\theta$}
\end{align*}

Rectangular coordinates of $(r,\theta+\pi)$:
\begin{align*}
x_2 &= r\cos(\theta+\pi) = r \cdot \ans{$(-\cos\theta)$} = \ans{$-r\cos\theta$} \\
y_2 &= r\sin(\theta+\pi) = r \cdot \ans{$(-\sin\theta)$} = \ans{$-r\sin\theta$}
\end{align*}

Therefore $x_1 = x_2$ and $y_1 = y_2$, so the two representations name the
\ans{same} rectangular point. \qquad $\square$

\vspace{6pt}
(c)\quad \ans{This shows that polar representations are not unique --- two
formally different pairs can describe the same location, unlike rectangular
coordinates where each point has exactly one $(x,y)$ description.}

\end{tcolorbox}

\end{document}
